% 【专题研讨】所属：第4章（由 chapters/revised/04-sources.tex \input 引入）
\minitopic{专题研讨：知觉证据、专长与认知渗透}“我亲眼看见”常被当作争论终点，哲学分析却要把这句话拆开。主体可能有某种视觉经验，经验可能由对象正常造成，主体可能把经验当作理由，也可能具备在相关环境中识别对象的能力。知觉证成理论的分歧，不只在于外部对象是否存在，而在于这些层面中哪一层提供认识资格，以及错误情形与成功情形共享什么。

\minitopic{经验作为理由还是通向事实的关系}教条主义认为，若主体有一个仿佛 \(p\) 的经验，且没有击败者，该经验即可给相信 \(p\) 初步确证。\parencite{pryor2000} 这种立场解释了为什么我们不必先证明自己不是缸中脑才相信面前有桌子。困难在于经验内容可能受偏见、期待或异常环境影响；若任何“看起来如此”都自动发放理由，错误经验似乎过于容易自我认证。 认识论析取论强调成功知觉与主观不可分辨的幻觉在认识地位上并不相同。真实看见时，世界中的事实本身能够进入主体的理由；幻觉最多给主体一个无责的表象。它因此拒绝“同样的内部经验必有同样的最高证成”这一假设。代价是要解释：主体在无法辨别两种情形时，怎样从第一人称角度管理自己的理由，以及事实作为理由是否仍需概念能力。 两种理论不应被误画成“经验论”对“无经验论”。析取论并不否认经验，教条主义也不必否认外部因果；争点是成功经验的认识价值能否完全由成功与失败共有的部分解释。设计案例时，需要固定主体的现象状态，再改变环境关系；也要固定环境关系，再改变主体是否把经验当作理由。只有一个维度变化，才能检验理论真正承诺。

\minitopic{认知渗透是否污染知觉}饥饿者觉得食物更醒目，球迷把模糊碰撞看成对方犯规，放射科医生从灰度图中看出新手看不出的病灶。若背景信念、欲望或情绪影响经验内容，称为认知渗透；若只影响注意、记忆、判断或报告，则未必是经验本身被渗透。心理机制难以从一句自我描述判断，因此哲学论证应与实验设计配合。 认识问题也不是“有影响就有污染”。专长驱动的影响可能提高敏感性，种族刻板印象驱动的影响可能制造系统偏差。一个候选规范是：影响若来自对真理有良好追踪记录、可被反馈纠正且对反例开放的能力，就可能提升证成；若影响由愿望直接绕过证据约束，则削弱证成。但可靠历史并非万能许可证，专家在分布改变、罕见疾病或自动化提示下仍会失准。 可以用四步诊断。第一，定位影响发生在刺激选择、早期加工、分类判断还是事后报告。第二，问影响是否对命题真值敏感：病灶特征消失时，专家是否撤回？第三，检查独立校验与反馈史。第四，分析主体是否拥有相关高阶证据，例如知道自己已看到带有暗示性的诊断标签。即使无法裁决经验内容是否改变，这四步仍能评价最终信念。

\minitopic{注意、忽视与“没有看见”}知觉能力不是被动接收全部环境信息。注意把有限资源分配给某些对象，也让其他对象留在背景。著名的非注意盲视类案例表明，刺激进入视野不等于主体获得关于它的可用证据。法律和医学判断中，“本可看见”必须进一步说明：对象是否清楚呈现、任务是否合理要求注意、当时负荷多大、主体是否受误导。 注意也有规范维度。驾驶员有义务监控行人和信号灯，却无须辨认每块广告牌；病理学家应系统扫描切片，而不是只寻找与初步诊断一致的区域。认识过失经常不是视觉器官失灵，而是查询策略失当。制度可以通过检查表、双人复核、随机化图像顺序和强制记录初步判断来改善注意结构。 但“应当注意”不能无限扩张。任何复杂场景都包含无数可见特征，事后知道结果后很容易把关键线索描述成显而易见。评价应采用事前标准：在不知道结局时，一个受过合理训练的主体会把该线索列入搜索范围吗？有无低成本方法能显著降低漏检？若没有，就不宜把所有知觉错误道德化。

\minitopic{多感官与仪器链}知觉证据常由多个通道共同产生。厨师用颜色、气味和质地判断食材；医生把影像、触诊和患者陈述结合；地质学家把卫星图、现场观察与化验结果对齐。通道一致可能提高信心，但前提是误差不过度相关。同一软件从一张原图生成三种可视化，不是三个独立知觉来源；三位医生都先看到同一算法标注，也可能共享锚定偏差。 仪器链应按转换步骤审计：对象如何产生信号，传感器如何采样，软件如何降噪或增强，显示器如何呈现，操作者如何分类。每一步都可能增加可见性，也可能删除信息。所谓“原始图像”往往已经经历曝光、重建和颜色映射。透明性不要求每位使用者掌握全部工程细节，却要求关键转换有记录、可回溯，并有与风险相称的校准。 把人—仪器系统看作扩展知觉者，有助于分配责任。设备制造者负责性能边界，机构负责校准和更新，操作者负责在适用范围内解释，决策者负责结合其他证据。若最终判断错误，不能只问“医生有没有看见”，还要问界面是否让异常可见、算法是否隐藏不确定性、工作流程是否允许复核。认识论由此进入设计层面。

\minitopic{专家之眼与环境变化}专家知觉并非把一个私密直觉贴上权威标签。它通常由大量有反馈的训练形成：专家知道看哪里、哪些差异稳定、何时请求第二意见。能力评价必须相对于任务与环境。鸟类学家在熟悉区域凭飞行轮廓识别物种，在迁徙路线改变后可能系统误判；放射科医生在本院设备上表现优良，换用不同成像协议时需要重新校准。 专长还会产生压缩：丰富线索被迅速组织成一个整体模式，专家不总能逐项报告。这使第一人称理由说明与实际可靠能力分离。制度不能因此放弃说明，而应使用适合专长的可审计方法，如要求标注关键区域、给出置信区间、保存初始读片并用盲评估测试表现。说明义务应帮助发现错误，而不是逼专家编造并未参与判断的事后规则。 新手与专家分歧也不自动由身份裁决。应问任务是否落在专长领域、专家是否掌握新手没有的信号、训练数据是否与当前环境匹配、是否存在利益冲突。专家权威是可撤销的认识信用，不是不可击败的社会等级。恰当信任专家与保留复核机制可以并存。

\minitopic{比较视野：蔽、心与察}荀子讨论“蔽”时，关注主体被局部遮蔽而不能把握更完整关系。\parencite{xunzi2014} 这可以与选择性注意、确认偏差和观点受限形成问题层面的比较：一个真实可见的局部如何因占据全部注意而产生总体误判？但“蔽”具有修身、论辩和政治秩序等更广背景，不能简单等同于某个现代心理学效应。 庄子通过视角转换展示事物在不同生命位置中的显现差异。\parencite{zhuangzi2009} 这不必导向“所有知觉同样正确”。更有生产力的问题是：何种差异来自对象关系的真实变化，何种差异来自判断标准，何种差异可通过移动位置或引入他者观点被校验？比较视野让知觉研究不只追问个体眼前有什么，也追问主体如何学习改变观看方式。

\minitopic{综合案例：辅助诊断的热区图}一套乳腺影像系统在图像上用红色热区提示模型关注位置。医生看到热区后发现细小阴影，判断需要活检；病理结果显示良性。事后审计发现模型热区并不稳定，同一图像轻微压缩后会移动，但医生指出的阴影确实存在且通常值得检查。这个案例不能用“模型错，所以医生没理由”或“医生看见了，所以模型无关”草率结束。 需要区分：热区是否改变了视觉经验、只改变注意，还是直接诱导分类；阴影本身给多少一阶支持；医生是否知道热区不稳定；活检决定的损益阈值如何；模型与医生错误是否独立。若无热区时医生通常漏掉该阴影，提示参与了能力实现；若热区随机指向也能引起同样判断，过程选择性不足。良性结果并不能反证当时活检决定不合理，因为决策目标是处理风险而非准确预言单个病理结果。 最后设计验证：保存无提示初读，随机显示真实、扰动和空白提示，测量定位、诊断、信心及反应时间；在不同设备与人群上复验；让医生说明哪些图像特征在移除提示后仍支持判断。这样的研究同时检验知觉证据、认知渗透、专家能力与制度可靠性。

\begin{tcolorbox}[breakable,colback=white,colframe=blue!30!black,title={专题研读任务},fonttitle=\bfseries]
选取一次“亲眼所见”的争议，画出从对象到信念的完整知觉链。标注刺激、注意、经验、分类、记忆、报告和行动各阶段，并为每个阶段设计一个只改变该环节的对照案例。
\end{tcolorbox}

\index{认知渗透}
\index{注意}
\index{专家知觉}
\index{仪器知觉}
\index{析取论}
